\section{留住每一次向前}
\subsection{自动保存与版本历史}
正常编辑会自动保存；离开前留意保存状态，网络断开时先保留本地文字。\menu{版本历史}按日期和时间展示项目的文本版本，可以查看差异、恢复内容，也可以让 AI 为改动生成摘要。摘要基于改动内容，不需要把整份作品重新讲一遍。

\menu{撤销}适合回退刚刚发生的一步，版本历史适合回到更早的时间点。两者都以实际保存的范围为准；二进制图片和附件建议另外导出留存。恢复前先看预览，避免把后来写好的内容一起换掉。

\subsection{把助手放在顺手的位置}
助手可以停靠在编辑器右侧、底部，也可以浮动。位置与大小可以按习惯调整。需要更宽的正文时，把助手放到独立标签页；它仍会跟随编辑器当前打开的文件和选区，也可以再放回编辑页。

选区小输入框与固定助手的输入草稿互不打扰，发送后仍使用同一段对话与上下文。对话标题处可以展开历史列表，新建对话不会删除旧讨论；需要接着写时，再选回原来的对话。

\subsection{一个轻一点的写作节奏}
先用一两句话说明你正在完成什么，给助手读真实材料的机会；让它先做一小段，再把有效的改法推广到相关位置。每次接受后看一眼正文和 PDF，需要停下来思考时，就只讨论。

这份手册可以留在书架上，也可以改名、删掉练习段落，变成你自己的第一份作品。工具尽量安静，写作仍在你手里。
